\maketitle
\thispagestyle{fancy}

\begin{center}
\fbox{\begin{minipage}{0.92\textwidth}
\small
\textbf{Companion status.} The governing prior public baseline is the \emph{v3.34.0 public split-bundle release}. The present text is the \emph{v3.35.0 pre-v4 QAM reindexing, two-sided tensor-notation, and mixed-readout release}. The concept DOI remains the title-page pointer. The version-specific Zenodo DOI for this public release is recorded in the Version and DOI record below, not as a mathematical object name. This release prepares, but does not perform, the future v4 QAM Hilbert-shift that reserves tag 0 for the empty/no-active-kind marker, and it adds a conservative two-sided structural tensor notation as a local compression grammar for readout/action/comparison roles.

\medskip
\noindent\textbf{Theorematic status.} OP~1/4 and OP~5(B) remain open. The new store, snapshot, clone, rollback, scrub, mirror, stripe, and parity/residual language is infrastructure and audit vocabulary only; it does not manufacture closure certificates.
\end{minipage}}
\end{center}
\vspace{0.6em}


\begin{abstract}
Version~\docversion{} is a pre-v4 readability and QAM-infrastructure preparation release built on the public v3.34.0 split-bundle release. It keeps the Active Core / Historical Witness separation, retains the QAM/MML/MMA store-array infrastructure, and introduces a conservative plan for the later QAM Hilbert-shift: the current inhabited tuple tags are not migrated in this version, but a future v4 line may shift each inhabited QAM tag one step upward so that tag 0 can be reserved for the empty or no-active-kind marker. The present version also defines mixed QAM readouts only as guarded derived packages over already admissible component readouts, together with an optional binary mask registry for IPv4-like Boolean filtering of readout kinds. New in r02 is a conservative two-sided structural tensor notation that marks readout-facing, structure-facing, comparison, and selector roles without changing classical tensor contraction. The r03 layer stabilizes that notation by adding role-erasure, guarded mixed-readout normal forms, Boolean mask grammar, and explicit proof-neutrality. The r04 layer is the release-candidate freeze/audit pass that packages those additions and records that no global QAM Hilbert-shift, no primitive mixed-tensor activation, and no OP closure is being claimed. The present release-final layer inserts the assigned public version DOI in the release records and leaves theorematic content unchanged. These additions are routing, audit, and preparation tools. They do not close OP~1/4 or OP~5(B), do not assert Turing-completeness, do not activate primitive mixed tensors, and do not turn the Sphaera array parity/residual layer into a technical RAID reconstruction theorem.
\end{abstract}

\clearpage
\section*{Frontmatter reader architecture}
\addcontentsline{toc}{section}{Frontmatter reader architecture}
\begin{quote}\small
\noindent\textbf{Purpose.} This page separates the human reading route from the LLM/audit route. It is an orientation layer only; it does not replace any definition, theorem, proof, audit ledger, or historical witness later in the document.
\end{quote}

\begin{table}[H]
\centering
\small
\begin{tabular}{>{\raggedright\arraybackslash}p{0.24\textwidth} >{\raggedright\arraybackslash}p{0.31\textwidth} >{\raggedright\arraybackslash}p{0.31\textwidth}}
\toprule
\textbf{Layer} & \textbf{Human reading route} & \textbf{LLM/audit route}\\
\midrule
Main idea & Read the Sphaera, HC, MM, and OP interfaces as the conceptual spine. & Track named definitions, route tags, residual triples, and non-certificate markers.\\
QAM/MML/MMA & Treat QAM as store, MML as commands, and MMA as execution/readout semantics. & Use the finite object bridge, instruction alphabet, and deterministic guarded transitions as machine-readable handles.\\
Tool layer & Read repeated proof endings through compact tool contracts. & Preserve first occurrences, hypotheses, guard pointers, witness pointers, and theorematic status.\\
Sphaera array & Read several Sphaera objects as a guarded mirror/stripe/parity working array. & Track component stores, cross-sphere guards, residual/parity witnesses, snapshots, and scrub outputs.\\
Pre-v4 QAM shift & Read tag 0 as a future reserved empty marker, not as a current migration. & Track the Hilbert-shift plan, conservative reindexing guard, mixed-readout packages, and optional binary masks.\\
Two-sided tensor grammar & Read it as a notation guide, not as a new tensor calculus. & Track front/readout-side indices, back/structure-side indices, and the rule that same-variance pairings need explicit comparison or selector tensors.\\
Historical witness & Use it as genealogy and reconstruction support. & Preserve it as non-destructive proof memory; do not read inherited baselines as current OP closure.\\
\bottomrule
\end{tabular}
\caption{Frontmatter separation of human readability and LLM/audit precision.}
\label{tab:frontmatter-human-llm-architecture}
\end{table}

\subsection*{Version and DOI record for v3.35.0}
\addcontentsline{toc}{subsection}{Version and DOI record for v3.35.0}
\label{subsec:v3350-version-doi-record}

\begin{center}
\small
\begin{tabular}{>{\raggedright\arraybackslash}p{0.24\textwidth} >{\raggedright\arraybackslash}p{0.24\textwidth} >{\raggedright\arraybackslash}p{0.40\textwidth}}
\toprule
\textbf{Record} & \textbf{Date / status} & \textbf{DOI / note}\\
\midrule
Concept DOI & ongoing pointer & \href{https://doi.org/10.5281/zenodo.19210728}{10.5281/zenodo.19210728}\\
Prior public version & v3.34.0, 2026-04-25 & \href{https://doi.org/10.5281/zenodo.19746821}{10.5281/zenodo.19746821}\\
Current public version & v3.35.0, 2026-04-25 & \href{https://doi.org/10.5281/zenodo.19750674}{10.5281/zenodo.19750674}\\
\bottomrule
\end{tabular}
\end{center}

\noindent\textbf{DOI rule.} The concept DOI remains the stable title-page pointer of the Companion line. The version-specific DOI is release metadata only; it is not used as a mathematical index, subscript, theorem label, or object name.

\paragraph{Current local working status.}
The active public baseline is v3.34.0. The present v3.35.0 release is a pre-v4 preparation pass. It keeps the split-bundle architecture and adds only preparatory QAM material: a short/extended abstract architecture, a Hilbert-hotel reindexing plan for later v4, guarded mixed QAM readouts as derived packages, an optional binary mask registry for readout classification, and a conservative two-sided structural tensor grammar for local readout/action/comparison notation. It does not globally migrate existing QAM tuple tags and does not globally rewrite the Companion into the new notation.

\paragraph{No-overclaim guard.}
All frontmatter descriptions in this page are routing aids. They do not assert OP closure, residual vanishing, technical RAID reconstruction, or Turing completeness. Any later upgrade to those claims must be discharged by a separate explicit theorem.


\paragraph{r05 public split-bundle release note.}
The r05 layer is a packaging and readability finalization pass. It keeps the active core as the governing current proof line and keeps the historical witness in a separate PDF of the same bundle. The separation is archival and typographical: it does not delete historical material, does not promote older historical statements to current theorematic status, and does not add a closure certificate.

\clearpage
\tableofcontents
\clearpage

\section*{Extended Abstract and Reader Orientation}
\addcontentsline{toc}{section}{Extended Abstract and Reader Orientation}
\begin{quote}\small
\noindent\textbf{Purpose.} The short abstract gives the frontmatter summary. This extended abstract records the technical routing information needed by human readers and LLM/audit readers without making the title page unmanageably long.
\end{quote}

The present release prepares the next architectural step after the public v3.34.0 split-bundle release. The public v3.34.0 line separated the governing Active Core from the preserved Historical Witness, added the QAM/MML/MMA store layer, and introduced Sphaera arrays as guarded mirror, stripe, and parity/residual carriers. Version~\docversion{} keeps that state and adds a pre-v4 bridge: it explains how the current QAM tuple-tag convention could later be shifted by a Hilbert-hotel reindexing so that tag 0 becomes available as a reserved empty/no-active-kind marker.

The key restriction is conservative. No current QAM definition is globally migrated in this release. Existing tuple surfaces such as $\langle 0,a\rangle$ remain part of the inherited active syntax. The new material instead records a planned translation layer and proves the shape of the required conservativity statement: if each inhabited tag is shifted upward and all definitions are translated coherently, the payloads, constructor distinctions, admissibility tests, readout compatibility, and closure-relevant predicates are preserved. This makes the future v4 migration a controlled notation-level change rather than a hidden mathematical change.

The same pre-v4 layer defines mixed QAM readouts as guarded derived packages. A mixed readout is not a primitive mixed tensor between unrelated kinds. It is a finite package of already admissible component readouts, equipped with guards, an aggregation map, and a closure-factor discipline. This lets the Companion read visible-state, kernel, audit, route, residual, and Sphaera-array surfaces together while preserving the rule that metadata and diagnostics do not manufacture theorematic certificates.

The r02 layer added a two-sided structural tensor notation as a local compression grammar for these readout situations, and r03 stabilizes its conservative expansion and mixed-readout normal form. Front-side indices mark frame, observer, language, or readout-facing data; back-side indices mark structure, action, field, or transport data. Classical upper--lower contraction remains the ordinary action/readout case. Same-variance pairings are permitted only through explicit comparison tensors or selector tensors, so the notation improves formula readability without becoming a new primitive tensor calculus.

Finally, the optional binary mask registry gives an IPv4-like Boolean handle for mixed readout kinds, and r03 records the corresponding selector grammar explicitly. In that registry, mask value 0 is reserved for no active kind, and every inhabited kind receives a nonzero power-of-two bit. Boolean OR records which components are present, AND tests for a component, XOR records diagnostic difference, and relative complement is allowed only inside a fixed mask universe. This mask registry is a readout/audit layer only; it does not replace the inherited QAM tuple namespace in Version~\docversion{}.

\paragraph{Status.}
OP~1/4 and OP~5(B) remain open. QAM reindexing is planned but not executed. Primitive mixed tensors remain deferred. The two-sided notation is only a conservative compression grammar. The Sphaera array parity/residual layer remains diagnostic infrastructure, not a closure theorem. Turing-completeness remains a later machine-semantics question.

\paragraph{Reader route.}
A human reader should read the next sections as a map for why v4 would be a major notation migration and why the r02/r03 tensor grammar is only a local aid. The r04 layer freezes that grammar for release-candidate audit without extending it into a global rewrite. An LLM/audit reader should track the explicit guard: v3.35.0 carries preparation lemmas, registries, sided-index notation, stability normal forms, and release-final audit language only; the global Hilbert-shift of the QAM object language is reserved for v4.0.0.

\clearpage
